\chapter{The dodecaphase closure of \texorpdfstring{\(\alpha\)}{alpha}}
\label{chap:alpha}
\index{alpha@\(\alpha\)}
\index{dodecaphase cycle}

This section establishes dodecaphase closure through integer arithmetic in
base one thousand. The construction makes explicit the twelve recurrences, the propagated
Euclidean quotients and the boundary conditions determining the output.

\section{The terminal selector of the three channels}
\label{sec:selector-terminal-alpha-rev6}

The closure brings together three already constructed generations, but does not identify them by
their decimal name. At terminal depth, the catalogs of compatible tails
have cardinalities
\[
 |\mathcal T_\pi|=196,\qquad
 |\mathcal T_e|=197,\qquad
 |\mathcal T_\varphi|=196.
\]
The selection object is therefore
\(\mathcal T_\pi\times\mathcal T_e\times\mathcal T_\varphi\), with
\[
 196\cdot197\cdot196=7\,567\,952
\]
possible triples.

The coordinate boundary fixed by the terminal reader is
\[
 q=(1,2,2,1,2,0)\in\mathbb F_3^6.
\]
The channel law \(\pi+e-\varphi\) has sign vector
\[
 \sigma=(1,1,-1)=(1,1,2)\in\mathbb F_3^3
\]
and nonzero dual line
\[
 \langle2\sigma\rangle
 =\langle(2,2,1)\rangle.
\]
The first datum acts on the boundary columns; the second, on
its rows. They are not two interchangeable features.

\begin{theorem}[Terminal matrix selector of \(\alpha\)]
\label{thm:selector-terminal-alpha-rev6}
In the preceding terminal catalog, the local signatures of Gram, norm, weight and
Hamming distance reduce the \(7\,567\,952\) triples to \(331\). Requiring
the column margin \(q\) leaves two candidates:
\[
 (120,197,157),\qquad(129,197,148).
\]
The condition that the row charge belong to the line of the channel
\(\langle\sigma\rangle\) selects exactly the second. Its blocks are
\[
 \boxed{
 b_\pi=021020,\qquad
 b_e=001220,\qquad
 b_\varphi=100210.}
\]
\end{theorem}

\begin{proof}
The two terminal matrices surviving the column margin are
\[
B_0=
\begin{pmatrix}
0&2&0&2&2&0\\
0&0&1&2&2&0\\
1&0&1&0&1&0
\end{pmatrix},
\qquad
B_1=
\begin{pmatrix}
0&2&1&0&2&0\\
0&0&1&2&2&0\\
1&0&0&2&1&0
\end{pmatrix}.
\]
Their row sums in \(\mathbb F_3\) are, respectively,
\[
 (0,2,0),\qquad(2,2,1).
\]
The first does not belong to
\(\langle(1,1,2)\rangle=\{0,(1,1,2),(2,2,1)\}\); the second does, since
\((2,2,1)=2\sigma\). Therefore only \(B_1\) simultaneously satisfies the
coordinate boundary and the channel charge. Its three rows are the blocks
shown and correspond to the triple \((129,197,148)\).
\end{proof}

The theorem does not obtain \(q\) from the row charge: \(q\) is the
coordinate datum of the terminal reader, and \(\sigma\) is the channel orientation.
What is proved is that, once both boundaries of the same object are declared,
the real branch is unique. This separation prevents a signature found
a posteriori from being turned into a supposedly derived law.

\begin{definition}[Two readings of the enriched terminal state]
Let
\[
 x_{\rm term}\in\mathcal X_{\rm TPK}^{\rm term,enr}
 \subseteq\mathcal X_{\rm TPK}^{[108]}.
\]
This state is the output
\[
 x_{\rm term}
 =(\mathcal U_{108}\circ\cdots\circ\mathcal U_1)(x_{\rm seed})
\]
of the complete pipeline of TPK selection, transport and update.
The preceding selector is the reading
\[
 \pi_{\rm term}(x_{\rm term})
 =B_{\rm term}
 :=\begin{pmatrix}021020\\001220\\100210\end{pmatrix}.
\]
The oriented register and the reader of the channels of \(90^\circ\) and
\(120^\circ\) produce, from the ledger of the same state,
\begin{equation}\label{eq:alpha-subespacio-firmado}
 \mathcal E_{108}^{90,120}\!
 \left(\mathcal R_{12}(x_{\rm term})\right)
 =(b^{90},b^{120},Q_{\rm TPK}),
\end{equation}
with
\[
\begin{aligned}
 b^{90}={}&(-495,-116,-684,108,601,-90,
             -173,-88,313,560,-397,461),\\
 b^{120}={}&(-425,-281,-481,671,-255,30,
              475,-543,680,251,6,-128),\\
 Q_{\rm TPK}={}&6263.
\end{aligned}
\]
For \(r=1,2,3\), let
\[
 B_r=\sum_{j=0}^{3}b^{120}_{r+3j},\qquad
 A_1=\frac{Q_{\rm TPK}+2B_1+B_2}{3},\qquad
 A_2=A_1-B_1,\qquad A_3=A_2-B_2,
\]
and \(d_{r,j}=b^{90}_{r+3j}\). The harmonic reader \(\Pi_H\) forms the
blocks
\begin{equation}\label{eq:alpha-modulo-imagen-H12}
 \Pi_H^{(r)}
 =\bigl(A_r,d_{r,1}-d_{r,3},
              d_{r,0}+d_{r,2},d_{r,0}-d_{r,2}\bigr)
\end{equation}
and interleaves them by columns. The signed reading is therefore the
produced composition
\begin{equation}\label{eq:alpha-Rsgn-definido}
 \boxed{
 R_{\rm sgn}:=
 \Pi_H\circ\mathcal E_{108}^{90,120}\circ\mathcal R_{12}:
 \mathcal X_{\rm TPK}^{\rm term,enr}
 \longrightarrow
 \mathcal U_{12}^{\rm sgn}:=\operatorname{im}R_{\rm sgn}.}
\end{equation}
Neither \(U\) nor \(K\) are coordinates added to the domain: \(U\) is the output
of \eqref{eq:alpha-Rsgn-definido}, and \(K\) will be its unique integral preimage
under Hadamard.
No factorization
\(B_{\rm term}\to U_{12}^{\rm sgn}\) is asserted: both outputs preserve distinct
components of the same terminal state.

For \(N'\geq N\geq108\), the truncations act only on the enriched
histories:
\begin{equation}\label{eq:alpha-truncamiento-firmado-preserva-KU}
 \tau_{N',N}:
 \mathcal X_{\rm TPK}^{[N'],\rm enr}
 \longrightarrow\mathcal X_{\rm TPK}^{[N],\rm enr}.
\end{equation}
If \(R_{\rm sgn,N}\) denotes the same composition at depth \(N\), naturality
is the concrete square
\begin{equation}\label{eq:alpha-cuadrado-naturalidad-Rsgn}
 \boxed{
 R_{\rm sgn,N}\circ\tau_{N',N}
 =R_{\rm sgn,N'}.}
\end{equation}
\end{definition}

\begin{theorem}[Reversible construction of the signed coordinate]
\label{thm:alpha-Rsgn-construido}
The composite \eqref{eq:alpha-Rsgn-definido}, applied to the produced terminal
state, gives exactly
\begin{equation}\label{eq:alpha-Rsgn-evaluacion}
 U=R_{\rm sgn}(x_{\rm term})
 =(2378,1406,2479,-452,998,-551,-668,-204,-371,-322,-28,-997).
\end{equation}
Its inversion by three Hadamard blocks of order four then produces
\[
 \boxed{
 K=(234,543,140,729,659,824,621,058,914,794,146,601).}
\]
The transformation is invertible; it satisfies the square
\eqref{eq:alpha-cuadrado-naturalidad-Rsgn} under the truncations
\eqref{eq:alpha-truncamiento-firmado-preserva-KU}; and it is independent of
\(\alpha\), CODATA and every metrological figure.
\end{theorem}

\begin{proof}
Formulas \eqref{eq:alpha-subespacio-firmado} and
\eqref{eq:alpha-modulo-imagen-H12} give
\[
 (B_1,B_2,B_3)=(972,-1073,101),\qquad
 (A_1,A_2,A_3)=(2378,1406,2479),
\]
and produce the three signed blocks
\[
\begin{aligned}
 U^{(1)}&=(2378,-452,-668,-322),\\
 U^{(2)}&=(1406,998,-204,-28),\\
 U^{(3)}&=(2479,-551,-371,-997).
\end{aligned}
\]
Their interleaving is \eqref{eq:alpha-Rsgn-evaluacion}. Since
\(H_4^2=4I_4\) and
\(\det(H_4\oplus H_4\oplus H_4)=(-16)^3=-4096\), we have
\begin{equation}\label{eq:alpha-Rsgn-inversa}
 K=\mathcal H_{12}^{-1}U
 =\frac14\mathcal H_{12}U
\end{equation}
after undoing the orbit permutation. Indeed,
\[
 H_4U^{(1)}=(936,2916,2484,3176),\quad
 H_4U^{(2)}=(2172,2636,232,584),
\]
and
\[
 H_4U^{(3)}=(560,3296,3656,2404).
\]
The forced division by four gives the orbits
\((234,729,621,794)\), \((543,659,058,146)\) and
\((140,824,914,601)\), whose interleaving is the declared seal.

As an independent check, the cyclic operators
\((D_3K)_m=K_m-K_{m+3}\) and
\((D_4K)_m=K_m-K_{m+4}\), together with \(Q(K)=\sum_mK_m\), satisfy
\begin{equation}\label{eq:alpha-extractor-90-120}
 \operatorname{rank}(D_3,D_4)=11,
 \qquad
 \operatorname{rank}(D_3,D_4,Q)=12.
\end{equation}
Thus \((D_3K,D_4K,Q(K))\) reconstructs a unique seal; moreover, the
calculation recovers exactly
\[
 D_3K=b^{90},\qquad D_4K=b^{120},\qquad Q(K)=Q_{\rm TPK},
\]
which certifies the prior reader of the register, without turning its outputs
into inputs. The truncation
\eqref{eq:alpha-truncamiento-firmado-preserva-KU} does not rewrite the first
one hundred and eight windows of the ledger; thus both sides of
\eqref{eq:alpha-cuadrado-naturalidad-Rsgn} read the same register and
produce \(U\). None of the matrices used contains \(\alpha\) or
an experimental datum.
\end{proof}

The observable 108-step calendar is a strictly poorer
projection: resetting its cursors every 9 steps erases precisely the
bidirectional coordinates that the enriched register preserves. Thus
the raw register does not replace the enriched state and is not used to
reconstruct \(R_{\rm sgn}\). The preceding selection theorem constructs
\(B_{\rm term}\); theorem \ref{thm:alpha-Rsgn-construido} constructs, separately,
the second publication of the same state.

\section{Dodecaphase state and reversible maps}

The theorem in this section receives from
theorem~\ref{thm:alpha-Rsgn-construido} the signed vector
\[
\Usgn=(2378,1406,2479,-452,998,-551,-668,-204,-371,-322,-28,-997).
\]
The typed route proved from that input is
\[
\Usgn\longleftrightarrow K,
\]
\[
K\longleftrightarrow(D_3K,D_4K,Q),
\]
\[
(P,E,\Phi,K,c_{13}=0)
\longleftrightarrow A=\alpha_{12}.
\]
Direct evaluation does not receive \(\alpha\) or a metrological value as
input. It receives \(\Usgn\), the three Archimedean representations and the
boundary condition. The
\cref{thm:selector-terminal-alpha-rev6} first constructs the selection of the
three terminal blocks; the present layer begins with the signed vector that
transports them to the dodecaphase closure. The reduced trace of a single
positional coordinate does not enter into this determination.

\begin{remark}[Two dodecaphase representations of different types]
\[
U_{12}^{\rm cat}=U_6\Vert U_6
\]
is a periodic word from the catalogue, whereas
\[
\Usgn=(2378,1406,2479,-452,998,-551,-668,-204,-371,-322,-28,-997)
\]
is a signed integer vector in the full state space of TPK\@. The first belongs to a
block representation and has period six. The second preserves orientation, is not
a base-thousand word and does not have that period. They are not interchangeable.
\end{remark}

\section{The canonical vector and the three orbits}

The dodecaphase vector is
\[
\boxed{
K=(234,543,140,729,659,824,621,058,914,794,146,601).}
\]
Its sum is
\[
Q(K)=6263.
\]
Step three divides \(\Cdoce\) into three orbits of four sectors:
\[
\begin{aligned}
K^{(1)}&=(K_1,K_4,K_7,K_{10})=(234,729,621,794),\\
K^{(2)}&=(K_2,K_5,K_8,K_{11})=(543,659,058,146),\\
K^{(3)}&=(K_3,K_6,K_9,K_{12})=(140,824,914,601).
\end{aligned}
\]
Each orbit traverses four sectors separated by \(90^\circ\). The
cardinal labels north--east--south--west appear after fixing origin and orientation; the
primary order is \((m,m+3,m+6,m+9)\).

\section{First representation: Hadamard transform}

Let
\[
H_4=
\begin{pmatrix}
1&1&1&1\\
1&1&-1&-1\\
1&-1&1&-1\\
1&-1&-1&1
\end{pmatrix}.
\]
The action on the three orbits gives
\[
\begin{aligned}
H_4K^{(1)}&=(2378,-452,-668,-322),\\
H_4K^{(2)}&=(1406,998,-204,-28),\\
H_4K^{(3)}&=(2479,-551,-371,-997).
\end{aligned}
\]
Interleaving by columns recovers
\[
\mathcal H_{12}(K)=\Usgn.
\]
Since
\[
H_4^2=4I_4,
\qquad
\det H_4=-16,
\]
the global transformation is invertible over \(\mathbb Q^{12}\):
\[
K^{(r)}=\frac14H_4U^{(r)},
\qquad
\det\mathcal H_{12}=(-16)^3=-4096\ne0.
\]
In the canonical representation, the inverse produces a unique integer vector in
\([0,999]^{12}\).

\begin{theorem}[Integral structure of the Hadamard transformation]
\label{thm:smith-hadamard}
The Smith normal form of the block \(H_4\), denoted by
\(\operatorname{SNF}(H_4)\), is
\[
 \operatorname{SNF}(H_4)=\operatorname{diag}(1,2,2,4).
\]
Consequently, for the direct sum of the three blocks,
\[
 \operatorname{coker}\mathcal H_{12}
 \cong(\mathbb Z/2\mathbb Z)^6\oplus(\mathbb Z/4\mathbb Z)^3,
\]
and the image of \(\mathbb Z^{12}\) has index \(4096\). A vector
\(u\in\mathbb Z^4\) belongs to \(H_4\mathbb Z^4\) exactly when
\[
 H_4u\equiv0\pmod 4.
\]
\end{theorem}

\begin{proof}
The greatest common divisor of the entries of \(H_4\) is one. The minors of
order two are even and one has absolute value two; the minors of order
three are multiples of four and one has absolute value four; finally,
\(\lvert\det H_4\rvert=16\). The determinantal divisors are
therefore \(1,2,4,16\), yielding the invariant factors
\(1,2,2,4\). The direct sum of three blocks repeats these factors and has
index \(16^3=4096\). Since \(H_4^{-1}=H_4/4\), the last
congruence is equivalent to the integrality of the preimage.
\end{proof}

\section{Relative cyclotomic transporter between sectors three and four}
\label{sec:transportador-ciclotomico-a34-rev8}

Let \(\mathsf R\) be the regular rotation of the cycle \(C_{12}\) on
\(\mathbb Q^{12}\), and let
\[
 W_d^{\rm cyc}:=\ker\Phi_d(\mathsf R)
\]
its cyclotomic sectors. The rational projectors onto the sectors of
orders three and four can be written
\[
 \Pi_3=
 \frac14(I+\mathsf R^3+\mathsf R^6+\mathsf R^9)
 -\frac1{12}\sum_{k=0}^{11}\mathsf R^k,
\]
\[
 \Pi_4=
 \frac16(I-\mathsf R^2+\mathsf R^4-\mathsf R^6
              +\mathsf R^8-\mathsf R^{10}).
\]
To fix the coordinate order, let
\[
 \mathsf P_{\rm orb}(z_1,\ldots,z_{12})
 =\bigl((z_1,z_4,z_7,z_{10}),
        (z_2,z_5,z_8,z_{11}),
        (z_3,z_6,z_9,z_{12})\bigr).
\]
The global transform used above is, in the natural order of \(C_{12}\),
\[
 \mathcal H_{12}
 :=\mathsf P_{\rm orb}^{-1}
   (H_4\oplus H_4\oplus H_4)\mathsf P_{\rm orb}.
\]
That is, \(H_4\oplus H_4\oplus H_4\) acts in the order of the three
step-three orbits, not on three contiguous quadruples in the natural order.
We then define the relative transporter
\[
 A_{34}:=\left.\Pi_4\mathcal H_{12}\Pi_3
          \right|_{W_3^{\rm cyc}}:
 W_3^{\rm cyc}\longrightarrow W_4^{\rm cyc}.
\]

\begin{theorem}[Dodecaphase cyclotomic transporter]
\label{thm:transportador-ciclotomico-a34-rev8}
In the bases
\[
 u_1=(1,-1,0)^{\times4},\qquad
 u_2=(0,1,-1)^{\times4}
\]
of \(W_3^{\rm cyc}\), and
\[
 v_1=(1,0,-1,0)^{\times3},\qquad
 v_2=(0,1,0,-1)^{\times3}
\]
of \(W_4^{\rm cyc}\), we have
\[
 A_{34}u_1=\frac23(v_1-v_2),\qquad
 A_{34}u_2=\frac23(v_1+v_2),
\]
and, therefore,
\[
 \boxed{\det A_{34}=\frac89}.
\]
In particular, \(A_{34}\) is a rational isomorphism between the two sectors.
\end{theorem}

\begin{proof}
The four preceding periodic sequences are substituted into the expressions
for \(\Pi_3,\Pi_4\) and into the three Hadamard blocks. The matrix of \(A_{34}\)
in the declared bases is
\[
 \begin{pmatrix}
  2/3&2/3\\
 -2/3&2/3
 \end{pmatrix},
\]
whose determinant is \(8/9\).
\end{proof}

The rational involution that exchanges \(W_3^{\rm cyc}\) and
\(W_4^{\rm cyc}\) by means of \(A_{34}\) and \(A_{34}^{-1}\), and fixes
\(W_{12}^{\rm cyc}\), transports the sum
\(W_4^{\rm cyc}\oplus W_{12}^{\rm cyc}\) to
\(W_3^{\rm cyc}\oplus W_{12}^{\rm cyc}\). In this latter sum, the cyclotomic
polynomial of the rotation is
\[
 \Phi_3(x)\Phi_{12}(x),
\]
which coincides with the polynomial of the Coxeter element of type \(E_6\). This is an exact
spectral signature and prepares the subsequent geometric corridor.

\begin{statusbox}[title={Rational transport and integral continuation}]
The operator \(A_{34}\) realizes rational cyclotomic transport between the
sectors of orders three and four. By fixing \(W_{12}^{\rm cyc}\), it places the
dodecaphase in the exact spectral sector
\(\Phi_3\Phi_{12}\) of type \(E_6\). The integral continuation is constructed
through two subsequent and compatible sources: the
\cref{thm:entrelazador-dodecafase-witt} provides the
\(M_{12}\)-equivariant isometry to the Witt module, and the
\cref{thm:identificacion-isometrica-tres-proyectores} chains the three
incidence projectors. On that integral realization, the
\cref{p12-83-c20-s6:thm:grafo-27-rectas-nuevo} constructs the graph of the
twenty-seven lines and the
\cref{p12-83-c20-s6:chap:holonomia-schlafli-e6-nuevo} identifies its Weyl
action. Thus, \(A_{34}\) fixes the rational stage of the genealogy and the
Witt--Schläfli intertwiners realize its lattice and
incidence extension.
\end{statusbox}

The sign separation is
\[
U_{12}^{+}
=(2378,1406,2479,0,998,0,0,0,0,0,0,0),
\]
\[
U_{12}^{-}
=(0,0,0,452,0,551,668,204,371,322,28,997),
\]
and
\[
\Usgn=U_{12}^{+}-U_{12}^{-}.
\]

\section{Second representation: transverse differences}

With indices in \(\{1,\ldots,12\}\) and the cyclic convention
\[
 K_{m+r}:=K_{1+((m+r-1)\bmod12)},
\]
we define
\[
(D_3K)_m=K_m-K_{m+3},
\qquad
(D_4K)_m=K_m-K_{m+4}.
\]
The dodecaphase transverse extractor is the operator
\[
 \mathcal E_{\mathrm{dod}}
 :=(D_3,D_4,Q).
\]
The subscript ``dod'' declares its domain. Shifts three and four
correspond to angular separations of \(90^\circ\) and \(120^\circ\) in the
oriented cycle \(C_{12}\), but \(\mathcal E_{\mathrm{dod}}\) is not the
hexad--octad incidence operator or the angular kernel subsequently used in
the towers.
The data are
\[
\begin{aligned}
D_3K={}&(-495,-116,-684,108,601,-90,\\
       &\qquad-173,-88,313,560,-397,461),\\[1mm]
D_4K={}&(-425,-281,-481,671,-255,30,\\
       &\qquad475,-543,680,251,6,-128),\\[1mm]
Q={}&6263.
\end{aligned}
\]
If \(D_3v=D_4v=0\), the vector is constant at steps three and four. Since
\(\gcd(3,4)=1\), those steps connect the entire cycle and
\[
\ker D_3\cap\ker D_4=\langle\mathbf1\rangle.
\]
Therefore
\[
\operatorname{rank}(D_3,D_4)=11.
\]
The total charge eliminates the constant translation, because
\(Q(\mathbf1)=12\ne0\). Therefore
\[
\boxed{\operatorname{rank}(D_3,D_4,Q)=12,}
\]
and
\[
K\longleftrightarrow(D_3K,D_4K,Q)
\]
is an exact coordinate system.

\begin{theorem}[Integral normal form of the transverse system]
\label{thm:smith-dodecafase}
Let
\[
 \mathcal D=(D_3,D_4,Q):
 \mathbb Z^{12}\longrightarrow\mathbb Z^{25}.
\]
Its nonzero invariant factors are
\[
 \operatorname{SNF}(\mathcal D)
 =\operatorname{diag}(\underbrace{1,\ldots,1}_{11},12),
\]
and, consequently,
\[
 \operatorname{coker}\mathcal D
 \cong\mathbb Z^{13}\oplus\mathbb Z/12\mathbb Z.
\]
\end{theorem}

\begin{proof}
The rows of \((D_3,D_4)\) are oriented incidences of the graph
\[
 \operatorname{Cay}
 \bigl(\mathbb Z/12\mathbb Z,\{\pm3,\pm4\}\bigr).
\]
Since steps three and four generate \(\mathbb Z/12\mathbb Z\), the graph
is connected.
The incidences of a spanning tree provide eleven unimodular reduced
minors; the first eleven invariant factors are therefore equal to
one.

Every minor of order twelve contains the row \(Q\), since the incidence
block has rank eleven. Adding the first eleven columns to the
last, a unimodular operation, makes the last column vanish in the incidence
rows and gives twelve in the row \(Q\). All maximal minors are
divisible by twelve. Choosing the eleven incidences of a tree gives a
remaining minor of absolute value one and, therefore, a maximal determinant
of absolute value twelve. The last invariant factor is exactly twelve.
\end{proof}

The component \(\mathbb Z/12\mathbb Z\) records the condition of
integral compatibility of the total charge with the reconstructed potential.
In particular, the rational uniqueness proved by the rank is refined here into
a statement about the integral lattice and its unique torsional obstruction.

The signed observable can also be reconstructed from these coordinates.
For \(r=1,2,3\), let
\[
B_r=\sum_{j=0}^{3}(D_4K)_{r+3j}.
\]
We obtain
\[
(B_1,B_2,B_3)=(972,-1073,101).
\]
Orbital sums are
\[
A_1=\frac{Q+2B_1+B_2}{3}=2378,
\quad
A_2=A_1-B_1=1406,
\quad
A_3=A_2-B_2=2479.
\]
If
\(d_{r,j}=(D_3K)_{r+3j}\) for \(j=0,1,2,3\), with the same cyclic
convention, the signed block of the orbit \(r\) is
\[
(A_r,d_{r,1}-d_{r,3},d_{r,0}+d_{r,2},d_{r,0}-d_{r,2}),
\]
exactly the corresponding Hadamard block.

\begin{remark}[Codomains of \(\mathcal H_{12}\) and
\(\mathcal E_{\mathrm{dod}}\)]
The correct notation is
\[
\mathcal H_{12}(K)=\Usgn,
\qquad
\mathcal E_{\mathrm{dod}}(K)=(D_3K,D_4K,Q).
\]
Both representations are compatible because they reconstruct the same vector,
but no diagram between their codomains has been defined here. Their formulas and
codomains are distinct. Sharing the steps \(3/4\) or the indices \(90/120\) does not
authorize identifying them.
\end{remark}

\section{The three Archimedean coordinates}

The vectors of twelve triads are
\[
\begin{aligned}
P={}&(141,592,653,589,793,238,462,643,383,279,502,884),\\
E={}&(718,281,828,459,045,235,360,287,471,352,662,497),\\
\Phi={}&(618,033,988,749,894,848,204,586,834,365,638,117).
\end{aligned}
\]
In concatenated form:
\[
\begin{aligned}
\iota(P)&=141592653589793238462643383279502884,\\
\iota(E)&=718281828459045235360287471352662497,\\
\iota(\Phi)&=618033988749894848204586834365638117.
\end{aligned}
\]
\(P\) is the decimal coordinate shared by the five regions of
\(\pi\); the calculation does not assert that those regions are identical.

\section{Base-thousand carry recurrence}

We fix
\[
c_{13}=0
\]
and calculate from right to left, \(m=12,11,\ldots,1\):
\[
s_m=P_m+E_m-\Phi_m-K_m+c_{m+1},
\]
\[
A_m=s_m\bmod1000,
\qquad
c_m=\frac{s_m-A_m}{1000},
\qquad0\le A_m<1000.
\]
The table is printed in increasing index order, but each carry comes from the row
below.

\begingroup
\scriptsize
\setlength{\tabcolsep}{3.3pt}
\begin{longtable}{@{}r|rrrr|rr|rr@{}}
\caption{Complete dodecaphase closure in base thousand.}\label{tab:acarreo-alpha}\\
\toprule
\(m\)&\(P_m\)&\(E_m\)&\(\Phi_m\)&\(K_m\)&\(c_{m+1}\)&\(s_m\)&\(A_m\)&\(c_m\)\\
\midrule
\endfirsthead
\toprule
\(m\)&\(P_m\)&\(E_m\)&\(\Phi_m\)&\(K_m\)&\(c_{m+1}\)&\(s_m\)&\(A_m\)&\(c_m\)\\
\midrule
\endhead
1&141&718&618&234&0&7&007&0\\
2&592&281&033&543&0&297&297&0\\
3&653&828&988&140&\(-1\)&352&352&0\\
4&589&459&749&729&\(-1\)&\(-431\)&569&\(-1\)\\
5&793&045&894&659&\(-2\)&\(-717\)&283&\(-1\)\\
6&238&235&848&824&\(-1\)&\(-1200\)&800&\(-2\)\\
7&462&360&204&621&0&\(-3\)&997&\(-1\)\\
8&643&287&586&058&\(-1\)&285&285&0\\
9&383&471&834&914&\(-1\)&\(-895\)&105&\(-1\)\\
10&279&352&365&794&0&\(-528\)&472&\(-1\)\\
11&502&662&638&146&0&380&380&0\\
12&884&497&117&601&0&663&663&0\\
\bottomrule
\end{longtable}
\endgroup

The carries are
\[
(c_1,\ldots,c_{13})
=(0,0,0,-1,-1,-2,-1,0,-1,-1,0,0,0).
\]
Negative values of \(s_m\) are ordinary borrowings in base thousand. Euclidean
division uniquely fixes the remainder and quotient.

The output is
\[
\boxed{
A=(007,297,352,569,283,800,997,285,105,472,380,663),}
\]
and therefore
\[
\boxed{
\begin{aligned}
\alpha_{12}
&:=\frac{7297352569283800997285105472380663}{10^{36}},\\
&=0.007297352569283800997285105472380663.
\end{aligned}}
\]
This rational is the initial dodecaphase window of the compatible section
generated by TPK transport. If
\(\rho_{N',N}:C_{N'}^\alpha\to C_N^\alpha\) are its truncations, the
complete output is
\[
 \boxed{
 \alpha_{\mathrm{HMT}}
 :=\varprojlim_N\alpha_N,\qquad
 \rho_{N',N}(\alpha_{N'})=\alpha_N,\qquad
 \alpha_{36}=\alpha_{12}.}
\]
Phase return extends the memory and narrows the surviving
cylinders to arbitrary depth; it neither resets the state nor ends the
construction at thirty-six digits.

\section{The complete affine identity}

Coordinate by coordinate,
\[
\boxed{
P_m+E_m-\Phi_m-K_m+c_{m+1}
=A_m+1000c_m.}
\]
If
\[
C=(c_1,\ldots,c_{12}),
\qquad
C^+=(c_2,\ldots,c_{13}),
\]
the correct vector form is
\[
\boxed{P+E-\Phi-K+C^+-1000C=A.}
\]
To avoid identifying the finite tails with the complete constants,
we define
\[
 \tau_{36}(x)=10^{-36}\left\lfloor10^{36}
 \bigl(x-\lfloor x\rfloor\bigr)\right\rfloor,
 \qquad
 \kappa(K)=10^{-36}\iota(K).
\]
The exact decimal notation of the concatenated identity is then
\[
 \tau_{36}(\pi)+\tau_{36}(e)-\tau_{36}(\varphi)-\kappa(K)
 =\alpha_{12}.
\]
Without the explicit truncation, the embedding \(\kappa\) and the cocycle
\(C^+-1000C\), the coordinate equality would be false.

We define
\[
\iota(v_1,\ldots,v_{12})
=\sum_{m=1}^{12}v_m1000^{12-m}.
\]
Multiplying each recurrence by \(1000^{12-m}\) and summing makes the carries
telescope:
\[
\iota(P)+\iota(E)-\iota(\Phi)-\iota(K)-\iota(A)
=1000^{12}c_1-c_{13}.
\]
Since \(c_1=c_{13}=0\):
\[
\boxed{
\iota(P)+\iota(E)-\iota(\Phi)-\iota(K)=\iota(A).}
\]
The two remaining integers are
\[
\iota(K)=234543140729659824621058914794146601,
\]
\[
\iota(A)=007297352569283800997285105472380663.
\]

\section{Inversion and equivalence of the three representations}

The recurrence is inverted by
\[
K_m=P_m+E_m-\Phi_m+c_{m+1}-A_m-1000c_m,
\]
and the concatenated integer representation by
\[
\iota(K)=\iota(P)+\iota(E)-\iota(\Phi)-\iota(A).
\]
The complete diagram is
\[
\boxed{
K\longleftrightarrow\Usgn,
\qquad
K\longleftrightarrow(D_3K,D_4K,Q),
\qquad
K\longleftrightarrow A\mid(P,E,\Phi,c_{13}=0).}
\]
The first two are linear representations; the third is a reversible affine
map on the fiber
conditioned by \(P,E,\Phi\) and the boundary.

\begin{remark}[Commutativity of the closure]
The direct map obtains \(K\) from the signed representation and produces
\(A\). The inverse formula proves that the output and the boundary recover the
same vector without residual freedom. Invertibility does not constitute
independent statistical evidence; the commutativity of the three
representations is an algebraic property of the closure.
\end{remark}

Uniqueness is expressed at three levels of the same construction:
\begin{enumerate}
  \item with the signed vector, \(P,E,\Phi\) and the boundary fixed, the vector and the
  output are unique;
  \item the representations are reversible and commutative;
  \item the complete genealogy
  \(\mathrm{APP}\to\mathrm{TRIT}\to\mathrm{TPK}\) produces those primitives
  before the local closure, and its compatible truncations extend the
  identity to the coinductive section \(\alpha_{\mathrm{HMT}}\).
\end{enumerate}

\section{Dodecaphase boundary condition}

The first dodecaphase window uses \(c_{13}=0\) and ends at
\[
\ldots|380|663.
\]
This boundary is part of the domain of the reversible affine map. The
extension does not introduce another system: it iterates the same graded transport
of the TPK on
\[
 w_6\to w_{12}\to w_{18}\to w_{24}\to w_{30}\to R_{36}\to G_9
\]
and constructs the successive windows with phase return, memory advance and
compatible terminal carry. The
\cref{sec:l9s102:alpha-dos-vias} proves that the truncation squares
of this route and of the analytic reader commute and that their surviving
cylinders have singleton intersection.

\subsection{Integral class of the carry}

The terminal condition can be compared with the periodic problem. Let
\(S\) be the cyclic shift on \(\mathbb Z^{12}\) and, for an integer
base \(b\geq2\), define
\[
 \partial_b=S-bI.
\]
The periodic affine identity takes the form
\[
 A=P+E-\Phi-K+\partial_bC.
\]
For every \(h\in\mathbb Z^{12}\), the transformation
\[
 K\longmapsto K+\partial_bh,
 \qquad
 C\longmapsto C+h
\]
preserves \(A\). Thus, the choice of carries is a choice of
representative within an integral class.

\begin{proposition}[Periodic quotient of the carry]
\label{prop:cociente-periodico-acarreo}
On a cycle of length twelve,
\[
 \operatorname{coker}(S-bI)\cong
 \mathbb Z/(b^{12}-1)\mathbb Z.
\]
Concatenation in base \(b\), taken modulo \(b^{12}-1\), determines the
periodic class.
\end{proposition}

\begin{proof}
The cyclic module is identified with
\(\mathbb Z[t]/(t^{12}-1)\). Imposing the relation \(t=b\) produces
\(\mathbb Z/(b^{12}-1)\mathbb Z\). Evaluation of the coefficient
polynomial is concatenation in base \(b\), with orientation fixed
by the order of the twelve sectors.
\end{proof}

The closure of \(\alpha\) does not use the periodic quotient: it uses an oriented
chain with a right endpoint and requires \(c_{13}=0\). Euclidean division
then selects a unique representative. This difference explains
why the open boundary is part of the definition of the map and
cannot be replaced by a cyclic identification of the carries.

\section{Comparison with the metrological reference}

After generation, the integer arithmetic of the first window
verifies the recognition
\[
\boxed{
\alpha_{12}
=10^{-36}
\left\lfloor\frac{10^{36}}{137.035999084}\right\rfloor.}
\]
Its inverse is
\[
\alpha_{12}^{-1}
=137.03599908400000000000000000000001066588\ldots,
\]
and exactly reproduces the central string published in the 2018 adjustment by the
Committee on Data for Science and Technology (CODATA) \cite{Tiesinga2021CODATA}:
\[
\boxed{137.035999084.}
\]
This is not a floating-point coincidence. The thirty-six digits of
\(\alpha_{12}\) constitute the exact truncation of the reciprocal of the
printed central string. This finite check does not define
\(\alpha_{\mathrm{HMT}}\) or limit its depth: the coinductive section
already produced determines all its compatible windows, while CODATA only
subsequently recognizes the prefix that its experimental resolution allows
it to test.

Metrological recognition is not part of the recurrence. The mathematical
result is the exact identity of the two routes
\[
 \alpha_A=\alpha_B=\alpha_{\mathrm{HMT}},
\]
whose dodecaphase windows are determined by the signed vector, the
three generated words and the compatible boundaries. The CODATA table
only subsequently compares an already fixed output.

Comparison with subsequent metrological adjustments is a separate external
check. The HMT output remains frozen and can therefore be tested without
modifying the generator.

\section{Coordinate independence and rigidity}

The closure tests distinguish deep modifications from local
perturbations:
\begin{center}
\small
\begin{tabularx}{\textwidth}{@{}lYY@{}}
\toprule
Variante & Result & What it controls \\
\midrule
canonical right--left protocol & igualdad exacta & reference \\
use \(U_6\Vert U_6\) as \(K\) & fails macroscopically & type distinction \\
rotar \(K\) & breaks the closure & dodecaphase origin \\
reflejar \(K\) & reverses the boundary & orientation \\
left--right carry & changes the word & direction of transport \\
perturb \(K_i\) by one unit & the exact word fails & rigidez profunda \\
\bottomrule
\end{tabularx}
\end{center}
A local perturbation may alter only a deep digit; it remains a
refutation of the exact equality. Structural variants produce larger
failures and show that the closure is not invariant under arbitrary orders.

\begin{remark}[Structure of the dodecaphase closure]
The closure of \(\alpha\) comprises an oriented cycle, two exact
representations of the vector \(K\), twelve typed inputs, a right-to-left
recurrence, thirteen carries, a boundary condition and a telescoping integer
identity. These data determine the output reproducibly.
\end{remark}
